\section{Unrestricted Problems and Fairness}\label{sec:unrestricted}

For partially ordered problems without restrictions on the hierarchy, plan existence is
undecidable already for deterministic problems \parencite{erol:et:al:amai-96}. For the fixed-method
concepts and for weak and strong method-based policies it is still semi-decidable: a solution has a
finite witness (an achievable primitive network with a solvability tree, a run, or a finite winning
strategy of the game of Theorem~\ref{thm:mp-andor}), and these witnesses can be enumerated and
checked. \CBb{} state that ``the general case is undecidable'' and do not distinguish the strong
cyclic column. We show that fair strong cyclic plan existence is of a different nature.

\begin{theorem}\label{thm:sigma11}
  Fair strong cyclic method-based plan existence for unrestricted problems is
  \begin{enumerate}[label=(\roman*)]
    \item not recursively enumerable,
    \item not co-recursively enumerable,
    \item in $\Sigma^1_1$, and
    \item $\Sigma^1_1$-hard.
  \end{enumerate}
  Hence it is $\Sigma^1_1$-complete and not arithmetical.
\end{theorem}

Items (i)--(iii) are fully formalised. Item~(iv) is formalised except for Harel's theorem, see
below. Note that (iv) implies (i) and (ii); we nevertheless prove them separately because
their proofs are elementary and self-contained.

\subsection{Fair solutions as infinite runs}

The proofs of (i) and (iv) use one construction. Let $M$ be a two-stack machine with initial
stacks $u,v$. We build a FOND problem whose task network holds the two stacks as two unordered
chains of compound tasks, one task per stack symbol, each chain closed by a task for its bottom.
The state holds the control state of $M$.
\begin{itemize}
  \item For each transition of $M$ that pops $x$ from stack $b$, the task for $x$ on stack $b$ has a
    method that executes an action $a_t$ for the transition and then pushes the new symbols, i.e.\
    replaces $x$ by the tasks for the pushed word. The action $a_t$ requires the source control state
    and has two effects: the \emph{continue} effect moves to the target control state, and the
    \emph{tick} effect adds a fact $\stopf$. Every transition action requires that $\stopf$ is
    false.
  \item Every stack symbol and every bottom also has a \emph{drain} method, refining it into an
    action that requires $\stopf$ and does nothing.
\end{itemize}
Since a symbol task can only be refined at the top of its stack, a policy that does not drain
simulates $M$ step by step. The network becomes empty, which is the only way to reach a goal, only
after a tick has added $\stopf$ and everything has been drained.

\begin{lemma}\label{lem:fair-inf}
  The problem has a fair strong cyclic solution iff $M$ has an infinite run from the initial
  configuration.
\end{lemma}

\begin{proof}
  Suppose $M$ has an infinite run. The policy simulates it. From every node of its execution
  structure a goal is reachable: either $\stopf$ already holds and the policy drains, or the next
  transition of the run is about to be executed, and its tick outcome leads to draining. Conversely,
  let a fair policy be given. As long as $\stopf$ is false the policy cannot drain, so at every
  non-goal node it must execute a transition (no node may be a dead end, since a goal must remain
  reachable). Following the continue outcomes gives an infinite run of $M$.
\end{proof}

\begin{proof}[Proof of Theorem~\ref{thm:sigma11}(i)]
  Consider a deterministic two-stack machine that simulates a universal Turing machine on input
  $n$, as used in the undecidability proof of deterministic HTN planning. It has an infinite run iff
  the $n$-th partial recursive function diverges on $0$. Divergence is not recursively enumerable,
  and the map from $n$ to the problem of Lemma~\ref{lem:fair-inf} is computable, so fair plan
  existence is not recursively enumerable.
\end{proof}

\begin{proof}[Proof of Theorem~\ref{thm:sigma11}(ii)]
  Halting reduces to deterministic plan existence for partially ordered HTN problems, via the
  intersection of context-free languages \parencite{erol:et:al:amai-96}. A deterministic problem
  has a fair strong cyclic solution iff it has a plan (Proposition~\ref{prop:fair-basic}(iii)), and
  the goal can be added as a final task. So halting reduces to fair plan existence. Since
  non-halting is not recursively enumerable, neither is the complement of fair plan existence.
\end{proof}

\subsection{Recurrence and $\Sigma^1_1$-hardness}

We now refine Lemma~\ref{lem:fair-inf}. Fix a control state $q$ of $M$ and keep the tick effect only
for the transitions that \emph{enter} $q$. A run of $M$ is \emph{$q$-recurrent} if it enters $q$
infinitely often.

\begin{lemma}\label{lem:fair-rec}
  The modified problem has a fair strong cyclic solution iff $M$ has a $q$-recurrent run from the
  initial configuration.
\end{lemma}

\begin{proof}
  If $M$ has a $q$-recurrent run, the policy simulates it. From any node, the run enters $q$
  again after finitely many steps, and the tick of that transition leads to a goal. So a goal is
  reachable from every node, by induction on the distance to the next entry into $q$.

  Conversely, let a fair policy be given and consider its attractor ranks
  (Theorem~\ref{thm:mp-fair-game}). At every node the policy has a move that witnesses the rank:
  all outcomes stay among the policy's nodes, and one outcome has a smaller rank. As before, this
  move executes a transition. If the witnessing outcome is the tick, the transition enters $q$.
  Otherwise the continue outcome has smaller rank. Following the witnessing moves, and the continue
  outcome whenever the witness is the tick, we obtain an infinite run in which every step either
  enters $q$ or decreases the rank. Since ranks are natural numbers, $q$ is entered infinitely
  often.
\end{proof}

Note that the proof uses the liveness part of fairness (the attractor). The safety part alone,
that the policy never gets stuck, only yields an infinite run, as in Lemma~\ref{lem:fair-inf}.

\begin{proof}[Proof of Theorem~\ref{thm:sigma11}(iv)]
  \textcite{harel:jacm-86} showed that the recurrence problem for nondeterministic Turing machines
  (does the machine, started on input $w$, have a run that visits a given state infinitely often?)
  is $\Sigma^1_1$-complete. This holds already for a fixed machine, since Harel's reduction is
  computable and a universal machine can compute and simulate its output. A nondeterministic
  two-stack machine simulates a Turing machine with finitely many steps per step and enters a
  designated state exactly when the simulated machine does. Hence recurrence for a fixed two-stack
  machine, with the input on one stack, is $\Sigma^1_1$-complete. By Lemma~\ref{lem:fair-rec} it
  reduces to fair plan existence.
\end{proof}

The formal development contains Lemma~\ref{lem:fair-rec} and the computable reduction. Harel's
theorem and the folklore simulation in the last proof are not formalised, since the
analytical hierarchy is not available in Mathlib.

\subsection{Membership in $\Sigma^1_1$}

\begin{proof}[Proof of Theorem~\ref{thm:sigma11}(iii)]
  By Theorem~\ref{thm:mp-fair-game}, a problem has a fair solution iff the initial node belongs to
  a fair-closed set of nodes. We play the game on finite descriptions of nodes (a state and a
  task network written with explicit identifiers), with progression only, so every configuration
  has finitely many successors, which can be computed. A fair-closed set together with its attractor
  can be described by two functions on configurations: membership $x(c)\in\{0,1\}$ in the set, and
  an attractor rank $r(c)$, where $0$ means ``not in the attractor'' and $k+1$ means ``in the
  attractor with rank $k$''. They have to satisfy a \emph{local} condition at every configuration
  $c$:
  \begin{itemize}
    \item if $x(c)=1$ then $r(c)>0$;
    \item if $r(c)=k+1$, then $c$ is a goal; or the planner has a move with a successor $c'$ such that
      $x(c')=1$ and $0<r(c')\le k$ (for decompositions); or the environment moves at $c$, all
      successors $c''$ have $x(c'')=1$, and some successor $c'$ has $0<r(c')\le k$ (for primitive
      tasks).
  \end{itemize}
  Conversely, the least ranks of a fair-closed set satisfy these conditions. Encoding
  configurations by natural numbers, the pair $(x,r)$ becomes a single function
  $f:\mathbb N\to\mathbb N$. A finite prefix of $f$ determines $x$ and $r$ for all configurations
  with small codes, and the relation $R(w,f(0)\cdots f(n-1))$ checks the initial configuration and
  the local condition at every configuration whose code and whose successors' codes are below $n$.
  This relation is computable (in fact in polynomial time), and the problem described by $w$ is
  solvable iff some $f$ passes $R$ on all of its prefixes.
\end{proof}

Note the contrast with the strong case. There the witness, a finite strategy, can be found by
enumeration. Here the witness must certify that a goal \emph{remains} reachable along unboundedly
growing networks, and no finite object can do that in general.
